\section{The AEP Protocol}
\label{sec:protocol}

\subsection{Ordering}

The runner's sequence, with the property each step exists to serve:

\begin{enumerate}
\item Acquire the lease with jittered retry.
\item Require that execution state already exists.
\item Prepare, vault and read back the request \emph{before any intent
      exists}, so that the intent, once written, names a request that is
      already durable.
\item Write the intent \textsc{none}$\rightarrow$\textsc{about-to-fire} by
      fenced CAS, \textbf{and} run the durability barrier on the \emph{same
      pinned connection}.
\item Convert the barrier's acknowledgement into a Redis-visible dispatch
      authorization.
\item Run the atomic pre-dispatch preflight, which re-checks that
      authorization.
\item Mint a single-use dispatch capability.
\item Only then call the connector.
\end{enumerate}

Two of these steps can refuse, and both refusals happen \emph{before} any
provider bytes exist, so both may soundly record a confirmed failure. If the
barrier fails at step~4, the runner makes one fenced attempt to record
\textsc{failed-confirmed} with evidence \texttt{LOCAL\_NO\_DISPATCH}. If the
preflight rejects at step~6 --- a lost lease, a moved version, an absent
authorization --- the runner does the same, with a reason naming the preflight
rather than the barrier, provided it still owns the lease. Either way, if
ownership has already been lost the runner records nothing and leaves the
intent for recovery. Together these are the paths that produce the measured
result of \cref{sec:eval-prevention}; the outcome classifier keys on the
resulting status rather than on which of the two refused, so they are counted
alike.

\subsection{P1 --- Fenced state}
\label{sec:p1}

\begin{property}[Fenced state]
No committed state write can be superseded and then resurrected by a stale
writer. Every write to the state key is admitted only if, in one atomic Lua
invocation, (a)~the caller's ownership token equals the live value of the lock
key, and (b)~the stored version equals the caller's expected version and the
candidate's version equals expected${}+1$.
\end{property}

Both conjuncts are necessary and the failure of either is a known
pattern. Version-only CAS admits a writer that lost its lease but holds the
current version; token-only checking admits a writer rebasing on a stale
read. The pairing is the fencing-token argument made against lease-only
locking~\cite{kleppmann2016locking}, with the token check and the write placed
inside one script so they cannot interleave. The same script additionally
freezes, inside the atomic region, the execution identity, the schema version,
every envelope field outside a declared mutable set, every other ledger entry,
and the intent's immutable fields; transitions are append-only with a prefix
check, and ledger entries can never be deleted.

\paragraph{Declared residuals} A plain \texttt{SET} from any client with socket
access is unconstrained (there is no ACL). Expiry, not deletion, ends the
guarantee: after the retention floor elapses the record is gone and any writer
may create version~1 again. And --- confirmed by a probe rather than argued ---
an AOF rewind can un-fence: P1 is a statement about one monotonic Redis
timeline, and because lock keys are deliberately \emph{not} barriered, a lease
whose release was lost can reappear alongside a rewound version, so a
previously fenced writer satisfies both conjuncts. There is no local fix; the
fix is consensus, which is a non-claim (\cref{tab:nonclaims}).

\subsection{P2 --- Detectable ambiguity}
\label{sec:p2}

\begin{property}[Detectable ambiguity]
If a \emph{durable} \textsc{about-to-fire} intent exists, then under F1--F5 it
eventually reaches exactly one \emph{terminal automated} state ---
\textsc{fired-confirmed}, \textsc{failed-confirmed} or
\textsc{permanently-ambiguous} --- and no automated transition moves it
afterwards; it is never silently re-dispatched and never silently dropped.
This holds provided the recovery loop runs, survives its own pass, can
eventually obtain the lease, and is not made eligible before the dispatch
window closes.
\end{property}

\emph{Terminal automated} is exact rather than cautious. The only edges out of
\textsc{permanently-ambiguous} are the two operator transitions of
\cref{sec:p3}, so an intent that has been escalated can still be resolved
later --- by a person, holding the lease and the exact version, asserting
something the protocol cannot check.

The design question P2 raises is what makes ``barrier before dispatch'' more
than a convention. In this protocol it is a checked precondition, in four
links:

\begin{enumerate}
\item Through the supported API, the barrier issues a \texttt{DurabilityAck}
      only when \texttt{WAITAOF} returned at least one local fsync.
\item Ordinary callers cannot construct, subclass or copy
      \texttt{DurabilityAck}; the in-process registry detects reuse and binds
      the guard to \texttt{execution:intent:prepared\_version}, so a guard for
      one attempt cannot authorise another through the supported path.
\item Authorizing dispatch \emph{consumes} the ack, then runs a script that
      re-checks lease, version, status and binding digest and writes an
      authorization key.
\item The preflight script requires that exact value; an absent authorization
      defaults to a value that cannot match, so it fails closed.
\end{enumerate}

\paragraph{The irreducible gap, stated} Redis exposes no way for a Lua script
to verify that a previous \texttt{WAITAOF} succeeded. A complete in-store proof
is therefore impossible. Under trusted-code assumptions, an authorization
records that the supported path received and consumed its ordinary barrier
guard; it is not an independent proof that Redis fsynced. The HMAC and registry
are mutation/reuse checks, not a cryptographic security boundary. Python
underscore naming and closure placement do not prevent arbitrary same-process
code from calling the module-internal issuer without running the barrier. A
principal with direct Redis access can also write the authorization key, and a
caller inside the process can compose the binding service with a connector
directly, bypassing the runner. A stronger attacker model requires a separately
enforced process or service boundary, which this implementation does not have.

\paragraph{A fifth link, and it is a clock} The four links above run from the
barrier to the dispatch. They are silent about the other end, and P2 needs it:
recovery must not classify an intent whose worker may still be about to
transmit. Nothing structural prevents that --- both principals are acting on
the same fenced state, and a lease can expire at any moment. What prevents it
is a delay. An intent becomes eligible for recovery only after
$T_{\text{client}} + B + T_{\text{settle}}$ has elapsed since it was prepared,
which is sized to cover the whole dispatch window. Remove that delay and the
protocol loses an effect: recovery reads back \emph{not applied}, truthfully,
because the request has not been sent yet, and records a confirmed failure that
the dispatch then contradicts. \Cref{sec:modelchecking} produces this in ten
steps. The delay is therefore a precondition of P2's soundness rather than a
courtesy about late responses, and it is a real-time argument --- which is why
the residual named in \cref{sec:model-clocks} is the gap it cannot close.

\paragraph{Recovery classification} For an eligible intent, under a freshly
acquired lease and after re-validating status and version, the mapping from
observation to outcome is total (\cref{tab:classification}). Recovery has no
reference to the mutation call at all --- it can only read back --- and the
transition table has no edge back to \textsc{about-to-fire}, checked in Python
and re-checked inside the atomic script.

\begin{table}[t]
\centering
\caption{Recovery classification. The two rows in bold are where declared
ambiguity is produced by design rather than by failure, and they are what
\cref{sec:eval-rq1} measures.}
\label{tab:classification}
\small
\begin{tabular}{@{}p{0.50\columnwidth}p{0.40\columnwidth}@{}}
\toprule
Observation & Outcome \\
\midrule
capability undeclared / unrecognised & \textbf{ambiguous}, no query \\
\textsc{no-readback}                 & \textbf{ambiguous}, no query \\
applied (either querying class)      & confirmed \\
not-applied $+$ \textsc{authoritative} & refuted \\
not-applied $+$ \textsc{positive-only} & \textbf{ambiguous} (named violation) \\
conflict                             & \textbf{ambiguous} \\
unknown, budget remaining            & stay unresolved, backoff \\
unknown, budget exhausted            & \textbf{ambiguous} \\
\bottomrule
\end{tabular}
\end{table}

\paragraph{False-positive ambiguity is the price, and it is the right
direction} A crash after a durable intent but \emph{before} transmission yields
an intent that recovery must treat as ambiguous even though no effect
exists. This inflates the declared-ambiguity rate; it can never inflate the
undetected-duplicate rate. \Cref{sec:eval-rq1} shows this asymmetry directly in
the data: ambiguity is highest at exactly the crash point where no effect can
possibly exist.

\subsection{P3 --- Fail-closed liveness bound}
\label{sec:p3}

\begin{property}[Fail-closed bound]
Automated reconciliation of one intent terminates within a configurable attempt
budget and duration budget, after which the intent is placed in a terminal
automated state and the execution is withdrawn from normal scheduling.
\end{property}

Defaults are 8 attempts and 24 hours, with full-jitter backoff. Attempt
counting is monotone and re-checked inside the atomic script: an
unresolved-to-unresolved edge must increment the counter by exactly one. On
exhaustion the intent becomes \textsc{permanently-ambiguous}, the execution is
paused, no new intent may be created for it, and the only exit is an operator
transition that itself requires the lease and the exact version.

\paragraph{Declared residuals} The bound is consulted only on the
\emph{unknown} branch --- every other observation resolves immediately, so this
is not a gap in practice, but it is not a global watchdog either. Lease
contention does not consume budget. There is no escalation \emph{mechanism}:
``ejects to operator escalation'' means reaching a terminal state and pausing,
and nothing alerts. And escalated records still expire: a retention floor of 31
days is enforced for any ledger containing an unresolved or ambiguous record,
but a floor is not permanence, so after it elapses the escalation evidence is
deleted without an operator having acted. Both residuals are confirmed by
executable probes in the artifact, not asserted.

\subsection{The state machine}

\Cref{fig:statemachine} is the transition set exported from the implementation,
not a drawing of it, so the presence or absence of an edge is a fact about the
code. Three states are terminal for automated reconciliation:
\textsc{fired-confirmed}, \textsc{failed-confirmed} and
\textsc{permanently-ambiguous}, which is the declared-ambiguity corner. Nothing
returns to \textsc{about-to-fire}: once an intent leaves the dispatch window it
cannot re-enter it, and that absence is P2's ``never silently re-dispatched''.
The dashed edges are the operator transitions of \cref{sec:p3}, and they are the
only way out of that corner.

\begin{figure}[t]
\centering
% GENERATED by scripts/gen_state_machine.py -- do not edit.
% Edge set imported from aep_core.core.intents
% .LEGAL_INTENT_TRANSITIONS, so this figure cannot drift from
% the transition table the Lua script enforces.
% 10 edges, 6 states, 3 terminal for automated reconciliation.
\begin{tikzpicture}[
    every node/.style={font=\scriptsize},
    state/.style={draw,rounded corners=2pt,inner xsep=4pt,inner ysep=3pt,align=center,fill=black!3},
    terminal/.style={state,fill=black!8,very thick},
    elabel/.style={font=\scriptsize,inner sep=1pt,align=center},
    to/.style={-{Stealth[length=4pt]},shorten >=1pt,shorten <=1pt},
  ]
  \node[state] (none) at (0.0,0.0) {\textsc{none}};
  \node[state] (atf) at (0.0,-1.5) {\textsc{about-to-fire}};
  \node[state] (fu) at (0.0,-3.3) {\textsc{fired-}\\\textsc{unconfirmed}};
  \node[terminal] (fc) at (-3.08,-5.6) {\textsc{fired-}\\\textsc{confirmed}};
  \node[terminal] (flc) at (3.08,-5.6) {\textsc{failed-}\\\textsc{confirmed}};
  \node[terminal] (pa) at (0.0,-5.6) {\textsc{permanently-}\\\textsc{ambiguous}};

  \draw[to] (none) edge node[elabel,midway,right=2pt] {CAS $+$ barrier} (atf);
  \draw[to,out=180,in=90,out looseness=1.5,in looseness=0.45] (atf) edge node[elabel,pos=0.46,sloped,above=1pt] {declared success} (fc);
  \draw[to,out=0,in=90,out looseness=1.5,in looseness=0.45] (atf) edge node[elabel,pos=0.46,sloped,above=1pt] {declared failure} (flc);
  \draw[to] (atf) edge node[elabel,midway,right=2pt] {anything else} (fu);
  \draw[to,out=200,in=160,looseness=2.8] (fu) edge node[elabel,midway,left=2pt,align=right] {unknown;\\budget left} (fu);
  \draw[to,out=225,in=45] (fu) edge node[elabel,pos=0.63,sloped,above=2pt] {read-back applied} (fc);
  \draw[to,out=315,in=135] (fu) edge node[elabel,pos=0.58,sloped,above=2pt] {authoritative absence} (flc);
  \draw[to] (fu) edge node[elabel,pos=0.72,right=2pt] {budget spent} (pa);
  \draw[to,dashed,out=170,in=10] (pa) edge (fc);
  \draw[to,dashed,out=10,in=170] (pa) edge (flc);
\end{tikzpicture}

\caption{The intent state machine, generated from the transition table in
\texttt{aep\_core/core/intents.py}. Solid edges are automated; the dashed edges
from \textsc{ambiguous} are the operator-only transitions. There is no edge
back into \textsc{about-to-fire} --- that absence is what ``never silently
re-dispatched'' means, and it is enforced inside the atomic script rather than
by the caller.}
\label{fig:statemachine}
\end{figure}

\subsection{Model checking}
\label{sec:modelchecking}

We specify the protocol in TLA+ and check it with TLC. The specification is
transcribed from the implementation --- the Lua guards, the exhaustive
transition table, the creation fence, the recovery classification --- and
deliberately not from this section. A gate re-derives the transition table and
the status alphabet from the source and fails if the two disagree, because a
model checker asked about the wrong state machine does not complain: it reports
success.

The endpoint is modelled as \cref{sec:model-endpoint} describes it, and this is
the modelling decision the result rests on. It is never given an intent
identifier, so it cannot deduplicate and a second dispatch yields a second
effect; whether an effect is applied is chosen adversarially rather than
derived from anything the caller did; and ambiguous evidence is reachable on
every path, so no caller can ever separate \emph{no effect} from \emph{effect,
answer lost}. A god-view variable records what really happened at the endpoint,
and no protocol action may read it --- enforced textually across the
specification's protocol definitions, because a guard that could see it would
verify a protocol strictly stronger than the implementation while every
property still came back green.

\emph{Checked:} P1 as version monotonicity and the absence of status
regression; P2 as legal-edges-only, no re-entry into \textsc{about-to-fire}, at
most one unresolved intent, and eventual resolution under fairness; P3 as the
bounded attempt budget; and the trilemma as two invariants --- no undetected
duplicate, and no lost effect. Declared ambiguity is not a violation, and that
asymmetry is \cref{sec:trilemma} stated formally.

\begin{table}[t]
\centering
\caption{Seven of the nine configurations that must fail. Each removes one
stated assumption; \emph{steps} is the length of the counterexample TLC
returns. Two further configurations, not shown, are non-vacuity witnesses:
\textsc{p2}'s liveness property is an implication, and they establish that its
antecedent is reachable rather than letting it pass while meaning nothing.}
\label{tab:modelchecking}
\small
\begin{tabular}{@{}p{0.44\columnwidth}p{0.34\columnwidth}r@{}}
\toprule
Assumption removed & Property that breaks & Steps \\
\midrule
% Steps are counterexample depths from scripts/run_tlc.sh at the committed
% bounds; one row per file in formal/configs/. Depths are used rather than
% state counts because a failing run stops at its first counterexample, so its
% state count reflects the search order and is not reproducible.
A single Redis timeline   & P1 monotonicity      & 4 \\  % aof-rewind.cfg
A truthful \texttt{fsync} & no lost effect       & 8 \\  % write-loss.cfg
The barrier, plus a restart & no lost effect     & 8 \\  % b3-no-barrier-restart.cfg
A truthful endpoint       & no lost effect       & 9 \\  % untruthful-endpoint.cfg
\quad\emph{same, duplicate horn} & no undetected duplicate & 16 \\ % untruthful-endpoint-duplicate.cfg
The reconciliation delay  & no lost effect       & 10 \\ % no-reconcile-delay.cfg
Operator authentication   & no lost effect       & 12 \\ % operator.cfg
\bottomrule
\end{tabular}
\end{table}

\emph{Not checked, and the exclusions are the interesting half.} Every
cryptographic construction --- the request binding, the vault, the digest --- is
out of the model rather than modelled badly, because an uninterpreted function
standing in for a hash proves nothing about the real one and assuming it
unforgeable would assume exactly what is in question. The in-process
acknowledgement chain of the four links above is \emph{assumed}, not checked:
what the model checks is the part that is atomic and Redis-visible, the
authorization key and the preflight that requires it. Real time is not modelled;
where that cuts in the protocol's favour it is neutralised, since a lease may
expire at any instant, which is more adversarial than the budget, and where it
cuts against, it is the reconciliation delay named above. Liveness is checked
within the version bound, so it establishes that nothing gets stuck for a reason
internal to the protocol, not termination in an unbounded run.

\paragraph{Nine configurations exist to fail} Six configurations check that the
properties hold. The other nine each remove one stated assumption and must
produce the declared counterexample; a configuration that passes when it should
fail turns the build red, which is what keeps this from being a check that
cannot fail. \Cref{tab:modelchecking} is where the section's content actually
is. Two rows are worth reading twice. The AOF-rewind residual of
\cref{sec:p1} --- previously argued, then confirmed by a probe --- falls out in
four steps once the single-timeline assumption is dropped. And the
truthful-\texttt{fsync} row is the write-loss result of
\cref{sec:eval-durability} in the model: the barrier is issued, it reports
success, the write is discarded, and an effect ends up with no record at all.

\paragraph{The barrier is a query about durability, not a guarantee} The
ablated barrier of B3 returns success without asking anything. The real barrier
under write loss returns success because the server answers that it fsynced.
These are the same state in the model, and the specification expresses them with
one expression, which is the sharpest available statement of what the mechanism
is worth: \textbf{the barrier is not a durability guarantee but a question about
durability, and it inherits the truthfulness of whatever answers it}. The two
are not quite the same system --- the ablated barrier cannot report failure,
where the real one still fails when the server does not answer --- so write loss
removes the barrier's ability to inform while leaving its ability to fail.

\paragraph{What this establishes, and what it does not} TLC explored every
reachable state under these bounds and found no violation. That is a fact about
a specification of a few hundred lines, and four things bound what it transfers.
The model is \emph{not} the implementation: nothing connects them but a hand
transcription and one gate over the transition table and status alphabet, so
refinement is neither proven nor claimed, and a defect in the Lua's JSON
handling or arithmetic is invisible here. Redis is \emph{assumed}: script
atomicity, expiry and \texttt{WAITAOF} semantics are premises, and
\cref{sec:eval-durability} is the standing demonstration that one of them can be
false in a way nothing detects. The scope is \emph{small}: two workers, two
intents, five versions, one step, so a defect requiring a third worker or a
second step would not appear. And the properties not checked above are not
checked --- most of the implementation's security argument is among them. What
model checking changes is narrower than it looks and still worth having: the
composition of the fences, the state machine and the classification table was
previously an argument over code paths, and it is now a checked one --- together
with, in \cref{tab:modelchecking}, the price of each assumption it rests on.
